% answers_short.tex —— 2014 年全国Ⅰ卷（文科） · 简短答案速查
% 由 tools/make_short_answers.py 自 answers.tex 的 \gksolhead/\gkanswer 提取，
% 并与 2014/2014_数学真题_全国1文_答案要点.md 逐题比对后生成，勿手改（改源 MD 后重跑）。

\begin{gknotice}
\textbf{使用说明}：本节只给最终答案便于快速核对，完整推导见解析版。答案为编写者独立推导，\textbf{非官方标准答案}。
\end{gknotice}

\gkpart{选择题}
\begin{center}\begin{tabular}{ccccccccccccc}
\toprule
题号 & 1 & 2 & 3 & 4 & 5 & 6 & 7 & 8 & 9 & 10 & 11 & 12 \\
\midrule
答案 & B & A & B & A & A & B & C & B & D & C & B & A \\
\bottomrule
\end{tabular}\end{center}

\gkpart{填空题}
\begin{itemize}[leftmargin=2.6em]
  \item[\textbf{13.}] $\dfrac23$
  \item[\textbf{14.}] $A$ 城市
  \item[\textbf{15.}] $(-\infty,8]$
  \item[\textbf{16.}] $150$
\end{itemize}

\gkpart{解答题}
\begin{itemize}[leftmargin=2.6em]
  \item[\textbf{17.}] (1) $a_n=\dfrac{n+2}{2}$；(2) $2-\dfrac{n+4}{2^{n+1}}$
  \item[\textbf{18.}] (1) 各组频率/组距依次为 $0.006,0.026,0.038,0.022,0.008$；在题给空白直方图的五个组距内分别作对应高度的柱形；(2) 平均数估计为 $100$，方差估计为 $104$；(3) 不能认为符合；样本中不低于 $95$ 的频率为 $0.68<0.80$
  \item[\textbf{19.}] (1) 见解析版；(2) $\dfrac12$
  \item[\textbf{20.}] (1) $x^2+y^2-2x-6y+8=0$；(2) $l:y=2$，此时面积为 $4$；或 $l:2x+y-6=0$，此时面积为 $\dfrac{12}{5}$
  \item[\textbf{21.}] (1) $b=1$；(2) $(-1-\sqrt2,\sqrt2-1)\cup(1,+\infty)$
  \item[\textbf{22.}] (1) 见解析版；(2) 见解析版
  \item[\textbf{23.}] (1) $\begin{cases}x=2\cos\theta,\\y=3\sin\theta,\end{cases}\ (\theta\in[0,2\pi))$；$2x+y-6=0$；(2) 最大值 $\dfrac{22\sqrt5}{5}$，最小值 $\dfrac{2\sqrt5}{5}$
  \item[\textbf{24.}] (1) $4\sqrt2$；(2) 不存在
\end{itemize}
